\input{preamble}
\title{More than One Way to Skin a Cat:\\
Preserving Verified Modes in RLVR}
\author{Anonymous Authors}
\begin{document}
\maketitle

\begin{abstract}
Correctness alone does not reveal how many successful solutions a model can
produce. We introduce \mb{}, five executable domains that distinguish correct
outputs for the same problem. An inference-only hosted-model comparison finds
highly accurate yet concentrated Graph outputs. Replaying verified outputs
during Dr.GRPO and MaxRL improves held-out correctness and sampled breadth
across model scales and on harder tasks. These results distinguish discovering
a successful behavior from keeping it available.
\end{abstract}

\begin{figure}[!htbp]
\centering
\includegraphics[width=\linewidth]{figures/modecollapse_story.pdf}
\caption{\textbf{Correctness can rise while alternatives disappear.}
One graph admits six valid colorings. Bars pool four eight-sample draws from
matched Qwen2.5-3B runs. At pass 6.5, Dr.GRPO produces 22/32 correct from one
mode; ReplayMaxRL produces 32/32 across four.}
\label{fig:story}
\end{figure}

\section{Preserving alternatives for mathematical agents}
\label{sec:introduction}\label{sec:collapse}
Mathematical agents use tools and verifiers to search and revise
\citep{brown2024monkeys,snell2024testtime}, but can exploit alternative
computations only if the model proposes them. Reinforcement learning with
verifiable rewards (RLVR) gives familiar and newly discovered correct outputs
the same reward \citep{guo2025deepseekr1}. Accuracy can therefore improve while
verified outputs narrow (Figure~\ref{fig:story}).

\input{results/frontier_comparison_20260911_workshop_motivation.tex}

\section{\mb{} makes output identity executable}
\label{sec:modebench}
\href{https://huggingface.co/datasets/od2961/ModeBench\#level-1}{\mb{}} assigns each accepted output an exact, executable identity: a canonical
key whose distinct values define verified \emph{modes}
(Figure~\ref{fig:modebench-examples}). Keys capture computation routes in
Countdown and MathIR, tested program behaviors in Python, and valid
configurations in Graph and PantryPlan. Cosmetic aliases merge; distinct routes
can share an answer (key rules: Appendix~\ref{app:domain-overview}). This measures
task-defined output support; it does not establish diversity of latent reasoning strategies.

\begin{figure}[!htbp]
\centering
\includegraphics[trim={0 9bp 0 0},clip,width=\linewidth]{figures/modebench_examples.pdf}
\caption{\textbf{One prompt, different verified outputs.}
Keys retain a color vector, normalized computation tree, tested return vector,
algebraic state trajectory, or ingredient support. Cosmetic aliases merge;
these keys identify task-defined outputs rather than free-form reasoning.}
\label{fig:modebench-examples}
\end{figure}

\section{ReplayMaxRL: learn from new samples and remember successes}
\label{sec:method}\label{sec:maxrl-method}\label{sec:replay}
ReplayMaxRL combines MaxRL's sampling-aware learning from new samples
\citep{tajwar2026maxrl} with replay of earlier verified successes
(Figure~\ref{fig:verified-support-story}).

\begin{figure}[!htbp]
\centering
% Crop only blank PDF margins; keep the figure itself unchanged.
\includegraphics[trim={4bp 23bp 4bp 10bp},clip,width=\linewidth]{figures/verified_support_story.pdf}
\caption{\textbf{Discovery and preservation use the same verified samples.}
MaxRL learns from fresh rewards; replay retains exemplars by key and revisits
rare successes. Memory cannot protect a mode before it is discovered.}
\label{fig:verified-support-story}
\end{figure}

\begin{figure}[!htbp]
\centering
\includegraphics[width=\linewidth]{figures/experiment1_retention_comparator_matrix.pdf}
\caption{\textbf{Replay improves correctness and verified support across scales.}
(A) ReplayDr.GRPO minus Dr.GRPO. (B) Qwen2.5-0.5B alternatives and the initial
model minus Dr.GRPO. Daggers indicate four pairs; otherwise five. Averages use
shared seeds; intervals: Appendix~\ref{app:supporting}. Models: \href{https://huggingface.co/od2961/maxent-grpo-models/blob/main/experiments/E78/README.md\#models}{0.5B}, \href{https://huggingface.co/od2961/maxent-grpo-models/blob/main/experiments/E79/README.md\#models}{1B}, \href{https://huggingface.co/od2961/maxent-grpo-models/blob/main/experiments/E80-R1/README.md\#models}{3B}.}
\label{fig:cross-scale-terminal-effects}
\end{figure}

Bank $\mathcal B_x$ stores one policy-generated exemplar $e_c$ of token length
$|e_c|$ per admitted key for prompt $x$. Replay trains policy $\pi_\theta$ with
\begin{equation}
L_{\mathrm{replay}}(x)=-\frac{1}{|\mathcal B_x|}
\sum_{c\in\mathcal B_x}\frac{\log\pi_\theta(e_c\mid x)}{|e_c|}.
\label{eq:lm-verified-replay}
\end{equation}
Each retained key receives equal replay weight, independent of its frequency
in fresh rollouts. This gives rare discoveries a continuing training signal.
Replacing MaxRL with Dr.GRPO \citep{liu2025understanding},
which debiases GRPO \citep{shao2024deepseekmath}, yields ReplayDr.GRPO.
Appendix~\ref{app:algorithm} gives the full update and
Appendix~\ref{app:theory-replay} the idealized retention guarantee.

Related work on collapse \citep{cui2025entropy,dang2025diversity} and success
replay \citep{zhang2025rlep} appears in Appendix~\ref{sec:related}.

\section{Does memory preserve verified output support?}
\label{sec:experiments}\label{sec:results}
\texttt{pass@8} is the chance that eight samples contain any correct response;
\texttt{distinct@8} is their average number of distinct verified modes.
\emph{Extra modes}, \texttt{distinct@8}$-$\texttt{pass@8}, counts alternatives
beyond the first correct output. Both support measures remain coupled to correctness.

We train for eight passes with five registered seeds per comparison. Each domain
has 128 held-out prompts with four eight-sample draws; none enter the bank.
Dr.GRPO/MaxRL controls run the same bank code with zero replay gradient.
Figure~\ref{fig:cross-scale-terminal-effects} also compares reward normalization,
\href{https://huggingface.co/od2961/maxent-grpo-models/blob/main/EXPERIMENTS.md\#direct-comparators}{UCPO} \citep{lochab2026ucpo}, \href{https://huggingface.co/od2961/maxent-grpo-models/blob/main/EXPERIMENTS.md\#direct-comparators}{RLEP-Dr} success replay, and \href{https://huggingface.co/od2961/maxent-grpo-models/blob/main/EXPERIMENTS.md\#semantic-ablations}{semantic entropy}
(protocols: Appendix~\ref{app:experimental-design}).



\noindent\textbf{Replay improves held-out correctness and support.}
\label{sec:results-retention}
ReplayDr.GRPO improves average held-out correctness and verified support across
all five domains at each of the three model scales
(Figure~\ref{fig:cross-scale-terminal-effects}).
On Qwen2.5-0.5B Countdown, \texttt{pass@8} rises $.48\to.67$ and
\texttt{distinct@8} $.49\to1.64$. Uniform key balancing contributes to broader correct output support:
it adds $.32$ extra modes per eight samples over \href{https://huggingface.co/od2961/maxent-grpo-models/blob/main/experiments/E120-R1/README.md\#models}{frequency weighting},
averaged across five Qwen2.5-0.5B domains
(Appendix~\ref{app:frequency-replay-progress}).

\begin{figure}[!htbp]
\centering
\includegraphics[width=\linewidth]{figures/e118_all_scale_factorial_progress.pdf}
\caption{\textbf{MaxRL and replay at all three model scales.}
Tracks connect initial, trained, and replay models. Complete tracks average
within seed; Falcon Dr.GRPO uses four. Qwen2.5-3B MaxRL equally averages
five paired domain means ($n=5,3,5,3,1$): dashed, descriptive, no interval.
\href{https://huggingface.co/od2961/maxent-grpo-models/blob/main/experiments/E118/README.md\#models}{Archived models}.}
\label{fig:maxrl-factorial}
\end{figure}

\noindent\textbf{Replay improves both Dr.GRPO and MaxRL.}
\label{sec:results-maxrl}
Adding replay improves average correctness and verified support under both
objectives at Qwen2.5-0.5B and Falcon3-1B (Figure~\ref{fig:maxrl-factorial}).
At Qwen2.5-3B, the descriptive MaxRL/ReplayMaxRL means are $.55/.69$
correctness and $.76/1.11$ distinct modes. This five-domain summary uses
partial seed sets without an interval (Appendix~\ref{app:maxrl-scale-extensions}).

\begin{figure}[!htbp]
\centering
\includegraphics[width=\linewidth]{figures/modebench_level_admission.pdf}
\caption{\textbf{Replay across difficulty levels with Qwen2.5-0.5B-Instruct.}
Open/filled markers: Level 1/2. (A) Frozen-model success.
(B) Terminal means weight four complete domains equally within five paired
seeds, shared across both levels and all methods. Pantry arm counts remain
separate (Appendix~\ref{app:level2-interim}). \href{https://huggingface.co/od2961/maxent-grpo-models/blob/main/experiments/E119/README.md\#models}{Archived models}.}
\label{fig:level2-admission}
\end{figure}

\noindent\textbf{Replaying stored solutions also helps on harder problems.}
\label{sec:results-levels}
\href{https://huggingface.co/datasets/od2961/ModeBench\#level-2}{Level 2} preserves validators, keys, and training/evaluation support histograms while
increasing difficulty: Graph hides four vertices instead of three.
Four domains have complete five-seed factorials; PantryPlan is still running.
On Graph, Dr.GRPO/MaxRL reach $.20/.43$ \texttt{pass@8} and $.22/.55$
\texttt{distinct@8}; replay raises these to $.76/.74$ and $1.24/1.21$.
Countdown adds $.017$ extra modes under ReplayMaxRL; its correctness
interval includes zero (all terminal results: Appendix~\ref{app:current-campaign-results}).

\section{Conclusion and Limits}
\label{sec:conclusion}\label{sec:discussion}\label{sec:limitations}
Replay improves correctness and sampled breadth, but retains only discovered
modes. Fixed-bank scores still have a declining tail
(Appendix~\ref{app:fixed-bank-survival}). Exact mode probabilities, causal
accuracy mechanisms and benefits under changed constraints remain unverified.

\label{main:lastpage}
\clearpage
\bibliographystyle{plainnat}
\bibliography{example_paper}
\clearpage
\appendix
\input{appendix}
\end{document}
